\section{Objek Lanjutan: this, Constructor, dan Iterasi}

Keyword \code{this} di JavaScript merujuk ke objek yang ``memiliki'' konteks eksekusi saat ini. Dalam metode objek, \code{this} merujuk ke objek tersebut. Dalam fungsi biasa (non-strict), \code{this} merujuk ke \code{window}; di strict mode, \code{this} adalah \code{undefined}. Arrow function tidak memiliki \code{this} sendiri---ia mewarisi dari scope luar \cite{jsInfo}.

\begin{table}[htbp]
\centering
\begin{tabular}{|P{3.5cm}|P{4cm}|P{4.5cm}|}
\hline
\textbf{Konteks} & \textbf{Nilai this} & \textbf{Catatan} \\
\hline
Metode objek & Objek pemilik & \code{obj.method()} \\
\hline
Fungsi biasa & \code{window} / \code{undefined} & Strict mode: undefined \\
\hline
Arrow function & Dari scope luar & Lexical this \\
\hline
Constructor & Objek baru & \code{new MyObj()} \\
\hline
Event handler & Elemen target & \code{el.\allowbreak addEventListener(...)} \\
\hline
\end{tabular}
\caption{Nilai this dalam berbagai konteks JavaScript}
\end{table}

Fungsi \textbf{constructor} membuat objek baru dengan operator \code{new}. Konvensi: nama PascalCase. Di dalam constructor, \code{this} merujuk ke objek baru. Metode \code{Object.keys()}, \code{Object.values()}, dan \code{Object.entries()} memungkinkan iterasi objek; \code{for...in} mengiterasi key (termasuk inherited---gunakan \code{hasOwnProperty()} untuk filter) \cite{mdnJsGuide}.

\begin{konsep}
\textbf{Metode bind, call, dan apply.} Untuk mengontrol \code{this} secara eksplisit: \code{fn.call(obj, arg1, arg2)} memanggil \code{fn} dengan \code{this} = \code{obj}; \code{fn.apply(obj, [args])} sama tapi argumen dalam array; \code{fn.bind(obj)} mengembalikan fungsi baru dengan \code{this} terikat permanen. Berguna untuk callback yang kehilangan konteks \code{this} \cite{jsInfo}.
\end{konsep}

